% GENERATED FILE. Edit canonical lessons, then run npm run generate:books.
% canonical-source-hash: fnv1a64:d20ab787d123ca78
% canonical-lessons: IT-C493-cognato, IT-C493-cognata, IT-C493-bisnonno, IT-C493-parentela, IT-C493-assomigliare

\chapter{Parenti acquisiti}
\label{ch:it-parenti-acquisiti}
\hlchaptermodality{493}

\begin{chapteropening}
\textbf{By the end of this chapter:} \emph{I can describe the wider family: in-laws, great-grandparents and family likeness.}

Complete the last lesson of chapter 493: I can describe the wider family: in-laws, great-grandparents and family likeness.
\end{chapteropening}

\section[il cognato]{il cognato — brother-in-law}

\label{lesson:IT-C493-cognato}



\begin{quote}
Before the new one: say the Italian for to grow fond of, then the Italian for affectionate.
\end{quote}

\subsection*{You'll want to know: il cognato}
\textbf{il cognato} — \textquotedblleft{}brother-in-law\textquotedblright{}.

\begin{quote}
Mio cognato è il fratello di mia moglie.
\end{quote}

\begin{quote}
\emph{My brother-in-law is my wife's brother.}
\end{quote}

The husband of your sister, or the brother of your husband or wife.

\begin{cousinweb}
\emph{cognato} is inherited from Latin \emph{cognātus} \textquotedblleft{}related by birth\textquotedblright{} (EtymDB), from \emph{co-} \textquotedblleft{}with\textquotedblright{} and \emph{nātus} \textquotedblleft{}born\textquotedblright{}; in Italian it moved to relatives by marriage.
\end{cousinweb}

\begin{grammarlens}[title={What you've built}]
One more word for telling, explaining and describing at B1.
\end{grammarlens}

\subsection*{Your turn}
\begin{itemize}
\raggedright
  \item \emph{Say it:} \emph{il cognato}
  \item \emph{Say it:} \emph{il cognato}, once more
  \item \emph{Say it:} a sentence of your own with \emph{il cognato}
  \item \emph{Recall:} say the Italian for to grow fond of, then the Italian for affectionate, then say \emph{il cognato} again
\end{itemize}

\begin{tcolorbox}[breakable,colback=teal!4,colframe=teal!35!black,title={Before you move on}]
What does il cognato mean? (\textquotedblleft{}Brother-in-law\textquotedblright{}.) Say it once more.
\end{tcolorbox}
\section[la cognata]{la cognata — sister-in-law}

\label{lesson:IT-C493-cognata}



\begin{quote}
Before the new one: say the Italian for affectionate, then the Italian for brother-in-law.
\end{quote}

\subsection*{You'll want to know: la cognata}
\textbf{la cognata} — \textquotedblleft{}sister-in-law\textquotedblright{}.

\begin{quote}
Mia cognata viene a pranzo con i bambini.
\end{quote}

\begin{quote}
\emph{My sister-in-law is coming to lunch with the children.}
\end{quote}

The wife of your brother, or the sister of your husband or wife.

\begin{cousinweb}
\emph{cognata} is the feminine of \emph{cognato}, inherited from Latin \emph{cognātus} \textquotedblleft{}related by birth\textquotedblright{} (EtymDB), from \emph{co-} \textquotedblleft{}with\textquotedblright{} and \emph{nātus} \textquotedblleft{}born\textquotedblright{}.
\end{cousinweb}

\begin{grammarlens}[title={What you've built}]
One more word for telling, explaining and describing at B1.
\end{grammarlens}

\subsection*{Your turn}
\begin{itemize}
\raggedright
  \item \emph{Say it:} \emph{la cognata}
  \item \emph{Say it:} \emph{la cognata}, once more
  \item \emph{Say it:} a sentence of your own with \emph{la cognata}
  \item \emph{Recall:} say the Italian for affectionate, then the Italian for brother-in-law, then say \emph{la cognata} again
\end{itemize}

\begin{tcolorbox}[breakable,colback=teal!4,colframe=teal!35!black,title={Before you move on}]
What does la cognata mean? (\textquotedblleft{}Sister-in-law\textquotedblright{}.) Say it once more.
\end{tcolorbox}
\section[il bisnonno]{il bisnonno — great-grandfather}

\label{lesson:IT-C493-bisnonno}



\begin{quote}
Before the new one: say the Italian for brother-in-law, then the Italian for sister-in-law.
\end{quote}

\subsection*{You'll want to know: il bisnonno}
\textbf{il bisnonno} — \textquotedblleft{}great-grandfather\textquotedblright{}.

\begin{quote}
Il mio bisnonno è nato in campagna.
\end{quote}

\begin{quote}
\emph{My great-grandfather was born in the countryside.}
\end{quote}

\emph{La bisnonna} — great-grandmother; \emph{i bisnonni} — great-grandparents.

\begin{cousinweb}
\emph{bisnonno} is built from \emph{bis-} \textquotedblleft{}twice\textquotedblright{}, from Latin \emph{bis} (Lewis \& Short), and \emph{nonno} \textquotedblleft{}grandfather\textquotedblright{}, from Late Latin \emph{nonnus} \textquotedblleft{}a monk, a tutor\textquotedblright{} (Lewis \& Short).
\end{cousinweb}

\begin{grammarlens}[title={What you've built}]
One more word for telling, explaining and describing at B1.
\end{grammarlens}

\subsection*{Your turn}
\begin{itemize}
\raggedright
  \item \emph{Say it:} \emph{il bisnonno}
  \item \emph{Say it:} \emph{il bisnonno}, once more
  \item \emph{Say it:} a sentence of your own with \emph{il bisnonno}
  \item \emph{Recall:} say the Italian for brother-in-law, then the Italian for sister-in-law, then say \emph{il bisnonno} again
\end{itemize}

\begin{tcolorbox}[breakable,colback=teal!4,colframe=teal!35!black,title={Before you move on}]
What does il bisnonno mean? (\textquotedblleft{}Great-grandfather\textquotedblright{}.) Say it once more.
\end{tcolorbox}
\section[la parentela]{la parentela — relatives, kinship}

\label{lesson:IT-C493-parentela}



\begin{quote}
Before the new one: say the Italian for sister-in-law, then the Italian for great-grandfather.
\end{quote}

\subsection*{You'll want to know: la parentela}
\textbf{la parentela} — \textquotedblleft{}relatives, kinship\textquotedblright{}.

\begin{quote}
Tutta la parentela è venuta alla festa.
\end{quote}

\begin{quote}
\emph{All the relatives came to the party.}
\end{quote}

All your relatives together; \emph{grado di parentela} — degree of kinship.

\begin{cousinweb}
\emph{parentela} is built on \emph{parente} \textquotedblleft{}relative\textquotedblright{}, from Latin \emph{parēns, parentis} \textquotedblleft{}a parent\textquotedblright{} (Lewis \& Short), from \emph{parere} \textquotedblleft{}to give birth\textquotedblright{}.
\end{cousinweb}

\begin{grammarlens}[title={What you've built}]
One more word for telling, explaining and describing at B1.
\end{grammarlens}

\subsection*{Your turn}
\begin{itemize}
\raggedright
  \item \emph{Say it:} \emph{la parentela}
  \item \emph{Say it:} \emph{la parentela}, once more
  \item \emph{Say it:} a sentence of your own with \emph{la parentela}
  \item \emph{Recall:} say the Italian for sister-in-law, then the Italian for great-grandfather, then say \emph{la parentela} again
\end{itemize}

\begin{tcolorbox}[breakable,colback=teal!4,colframe=teal!35!black,title={Before you move on}]
What does la parentela mean? (\textquotedblleft{}Relatives, kinship\textquotedblright{}.) Say it once more.
\end{tcolorbox}
\section[assomigliare]{assomigliare — to look like, to resemble}

\label{lesson:IT-C493-assomigliare}



\begin{quote}
Before the new one: say the Italian for great-grandfather, then the Italian for relatives.
\end{quote}

\subsection*{You'll want to know: assomigliare}
\textbf{assomigliare} — \textquotedblleft{}to look like, to resemble\textquotedblright{}.

\begin{quote}
Assomigli molto a tuo padre.
\end{quote}

\begin{quote}
\emph{You look a lot like your father.}
\end{quote}

\emph{Assomigliare a qualcuno}; reflexive \emph{assomigliarsi} — to look alike.

\begin{cousinweb}
\emph{assomigliare} is built on \emph{somigliare} \textquotedblleft{}to resemble\textquotedblright{}, usually traced to Vulgar Latin \emph{similiāre}, from Latin \emph{similis} \textquotedblleft{}like\textquotedblright{} (Lewis \& Short).
\end{cousinweb}

\begin{grammarlens}[title={What you've built}]
One more word for telling, explaining and describing at B1.
\end{grammarlens}

\subsection*{Your turn}
\begin{itemize}
\raggedright
  \item \emph{Say it:} \emph{assomigliare}
  \item \emph{Say it:} \emph{assomigliare}, once more
  \item \emph{Say it:} a sentence of your own with \emph{assomigliare}
  \item \emph{Recall:} say the Italian for great-grandfather, then the Italian for relatives, then say \emph{assomigliare} again
\end{itemize}

\begin{tcolorbox}[breakable,colback=teal!4,colframe=teal!35!black,title={Before you move on}]
What does assomigliare mean? (\textquotedblleft{}To look like, to resemble\textquotedblright{}.) Say it once more.
\end{tcolorbox}
